% =====================================================================
\chapter{Résultats complémentaires}
\label{annexe-resultats}
% =====================================================================

Cette annexe rassemble les résultats que le \cref{chap-resultats} résume sans
les détailler~: le résultat de classement plus nuancé de l'évaluation~A
(\cref{sec-res_a_protocole,sec-res_a_classement}), le contrôleur budgété
(\cref{sec-res_ablation_c}), le système assemblé bout-en-bout
(\cref{sec-res_ablation_d}), les coûts de calcul observés
(\cref{sec-res_couts}), les limites du travail (\cref{sec-res_limites}) et les
directions qu'elles ouvrent (\cref{sec-res_futur}).

% =====================================================================
\section{Évaluation A, complément : protocole et classement}
% =====================================================================

\subsection{Deux contraintes de protocole}
\label{sec-res_a_protocole}

Deux caractéristiques du protocole de l'évaluation A conditionnent la lecture
des chiffres. \textbf{Un plafond de \num{25} tuiles DOM par page}
(\cref{sec-tuilage}) borne aussi l'évaluation~: aucune tuile au-delà n'a été
présentée au modèle générateur, donc aucune question ne peut en dépendre. Sur
\num{17} des \num{18} pages, le nombre de tuiles DOM est exactement \num{25}
(\cref{fig-res_unites_page})~: la valeur de \num{24.0} tuiles par page est
donc un \emph{plafond atteint}, non une mesure de densité de tuilage réelle,
qui serait supérieure. \textbf{Les trois conditions couvrent la même région de
page}~: le protocole calcule une frontière verticale --- le bas de la plus
basse tuile DOM retenue --- et y restreint la grille et les fenêtres de texte,
de sorte que les comparaisons de taux d'attachement et de budget de pixels
portent bien sur le même contenu.

\begin{figure}[t]
\centering
\includegraphics[width=0.9\textwidth]{res_unites_par_page.pdf}
\caption[Unités par page dans les trois conditions]{Nombre d'unités
récupérables produites par chaque condition, page par page, pour la même région
de page. Le plafond de \num{25} tuiles DOM est atteint sur \num{17} des \num{18}
pages~; seule la page la plus courte (\code{George\_P.\_Gunn}, \num{7} tuiles)
reste en dessous. La grille aveugle produit systématiquement entre \num{2} et
\num{7} bandes, soit un ordre de grandeur de moins.}
\label{fig-res_unites_page}
\end{figure}

\subsection{Résultat 3 : la qualité de classement}
\label{sec-res_a_classement}

\Cref{tab-res_ablation_a} rapporte le rang percentile (\cref{eq-pctrank}), sur
toutes les paires attribuables puis sur les \num{107} paires attribuables dans
les \emph{trois} conditions simultanément (comparaison strictement appariée).

\begin{table}[t]
\centering
\small
\caption[Évaluation A : tuilage, grille et texte plat]{Qualité de classement des
trois conditions, même scoreur TF-IDF, \num{18} pages et \num{409} questions.
\%rang désigne le rang percentile moyen~: \num{0.5} sous classement aléatoire,
invariant à la taille du vivier. Les Recall bruts ne sont \emph{pas}
comparables entre conditions (\cref{sec-eda_pool}) et sont donnés à titre de
diagnostic.}
\label{tab-res_ablation_a}
\begin{tabular}{@{}l r r r r r@{}}
\toprule
Condition & $n$ & unités/page & R@1 & R@5 & \%rang \\
\midrule
\multicolumn{6}{@{}l}{\emph{Toutes les paires attribuables}} \\
Tuiles DOM   & \num{212} & \num{24.0} & \num{0.547} & \num{0.774} & \num{0.868} \\
Grille aveugle & \num{108} & \num{4.9}  & \num{0.722} & \num{0.991} & \num{0.839} \\
Texte plat   & \num{232} & \num{12.3} & \num{0.612} & \num{0.875} & \num{0.867} \\
\midrule
\multicolumn{6}{@{}l}{\emph{Paires appariées (\num{107} attribuables dans les trois conditions)}} \\
Tuiles DOM   & \num{107} & \num{24.0} & \num{0.514} & \num{0.841} & \num{0.901} \\
Grille aveugle & \num{107} & \num{4.9}  & \num{0.729} & \num{0.991} & \num{0.845} \\
Texte plat   & \num{107} & \num{12.3} & \num{0.710} & \num{0.944} & \textbf{\num{0.928}} \\
\bottomrule
\end{tabular}
\end{table}

\begin{figure}[t]
\centering
\includegraphics[width=\textwidth]{res_ablation_a.pdf}
\caption[Évaluation A : attachement et classement]{\emph{À gauche}~: taux
d'attachement. \emph{À droite}~: rang percentile moyen sur toutes les paires
attribuables (barres pleines) et sur les \num{107} paires appariées (barres
hachurées). La ligne pointillée marque la valeur attendue sous classement
aléatoire.}
\label{fig-res_ablation_a}
\end{figure}

Trois lectures s'imposent. \textbf{L'artefact de vivier est massif}~: la
grille atteint \num{0.991} de Recall@5, un score presque parfait qui
n'exprime rien d'autre que le fait qu'elle ne propose que \num{4.9}
candidats par page --- ce qui valide empiriquement la nécessité du rang
percentile. \textbf{Le tuilage DOM domine la grille sur la métrique
corrigée}~: \num{0.868} contre \num{0.839} sur toutes les paires,
\num{0.901} contre \num{0.845} sur le sous-ensemble apparié --- un avantage
modeste (trois à six points) mais cohérent entre les deux ensembles, qui
s'ajoute au doublement du taux d'attachement et à la division par \num{2.7}
du budget de pixels. \textbf{Le texte plat devance le tuilage DOM sur les
paires appariées}~: \num{0.928} contre \num{0.901}, contrairement à la
prémisse partagée par PixelRAG, par MAGE-RAG et par notre propre
introduction.

\subsection{Pourquoi le texte plat reste compétitif}

Ce dernier résultat n'est pas un défaut de la condition textuelle~: trois
facteurs l'expliquent. \textbf{Le scoreur n'exploite aucune information
spatiale}~: TF-IDF, employé dans les trois conditions pour isoler l'effet du
découpage, ne sait pas qu'une cellule appartient à une ligne de tableau ou
qu'un titre coiffe un paragraphe --- il voit un sac de mots, précisément la
propriété sur laquelle repose la valeur du tuilage structuré. Le résultat
mesure donc le tuilage sur une métrique aveugle à ce qu'il apporte~; la
question de savoir si l'information spatiale préservée par des tuiles
alignées apporte un gain est traitée séparément par l'évaluation~B, qui
utilise un scoreur capable de la percevoir. \textbf{Le critère d'attribution
favorise le texte}~: une correspondance verbatim est structurellement plus
favorable à une condition qui manipule le texte directement qu'à une
condition qui relit l'aperçu textuel d'un recadrage d'image --- l'essentiel
de l'écart de taux d'attachement (\SI{56.7}{\percent} contre
\SI{51.8}{\percent}) en découle. \textbf{Le texte plat ne porte aucune
information visuelle}~: un chiffre lu sur un axe, une correspondance
ligne-colonne dans un tableau complexe, rien de cela ne subsiste dans du
texte extrait --- une limite que nos \num{18} pages, majoritairement
textuelles, ne font pas ressortir, mais qui se manifesterait sur un rapport
financier ou une notice technique. Cette évaluation établit donc que le
texte plat est \emph{compétitif sur ce corpus et sous ce scoreur}, non qu'il
est préférable dans l'absolu.

% =====================================================================
\section{Évaluation C : contrôleur budgété et sélection top-\texorpdfstring{$k$}{k} statique}
\label{sec-res_ablation_c}
% =====================================================================

\subsection{Protocole}

Cette évaluation mesure la capacité du contrôleur (\cref{sec-controleur}) à
rassembler des preuves réparties sur deux documents. Pour chacune des
\num{91} questions inter-pages construites depuis HotpotQA
(\cref{sec-protocole_multihop}), le graphe combiné des deux articles cités
est construit, avec recâblage des arêtes \code{links\_to} sur leur point
d'entrée réel (\cref{sec-graphe_corpus}). Le contrôleur et une référence
top-$k$ statique sont exécutés sur des scores de pertinence
\emph{identiques}, avec le même budget~; seule la politique de sélection
diffère. La métrique est la \textbf{couverture complète}~: les preuves
sélectionnées contiennent-elles au moins une tuile positive sur
\emph{chacune} des deux pages~? Deux configurations d'amorce sont
rapportées~: $k_{\text{seed}} = 3$ (valeur par défaut, budget \num{5}) et
$k_{\text{seed}} = 5$ (égale au budget).

\subsection{Diagnostic, correctif et résultat}
\label{sec-res_c_avant}

La première mesure, frontière propagée uniquement le long de
\code{reading\_order}, donnait \SI{29.7}{\percent} de couverture complète pour
le contrôleur contre \SI{40.7}{\percent} pour la référence à l'amorce par
défaut, et une égalité exacte à $k_{\text{seed}} = 5$ (bloc supérieur du \cref{tab-res_ablation_c}) --- attendue, puisqu'à amorce égale au budget le
contrôleur ouvre son amorce intégralement et \emph{devient} la référence
plutôt que d'y dégénérer, ce qui vérifie la cohérence de l'implémentation.
Le déficit à $k_{\text{seed}} = 3$ s'expliquait par l'incapacité de la
frontière à franchir une limite de document~: \code{reading\_order} ne relie
jamais deux pages, si bien que le second article n'était atteignable que si
l'amorce y avait déjà placé un nœud.

\begin{table}[t]
\centering
\small
\caption[Évaluation C : couverture inter-pages]{Fraction des \num{91} questions
inter-pages pour lesquelles les preuves sélectionnées couvrent une tuile
positive sur \emph{chacun} des deux articles, avant et après le correctif de
frontière inter-documents.}
\label{tab-res_ablation_c}
\begin{tabular}{@{}l r r@{}}
\toprule
Amorce $k_{\text{seed}}$ & Contrôleur budgété & Top-$k$ statique \\
\midrule
\multicolumn{3}{@{}l}{\emph{Avant correctif (\code{reading\_order} seul)}} \\
\num{3} (budget \num{5}) & \SI{29.7}{\percent} \enspace(\num{27}/\num{91}) & \SI{40.7}{\percent} \enspace(\num{37}/\num{91}) \\
\num{5} ($=$ budget)     & \SI{40.7}{\percent} \enspace(\num{37}/\num{91}) & \SI{40.7}{\percent} \enspace(\num{37}/\num{91}) \\
\midrule
\multicolumn{3}{@{}l}{\emph{Après correctif ($+$ \code{links\_to} et \code{continues})}} \\
\num{3} (budget \num{5}) & \SI{39.6}{\percent} \enspace(\num{36}/\num{91}) & \SI{40.7}{\percent} \enspace(\num{37}/\num{91}) \\
\num{5} ($=$ budget)     & \SI{40.7}{\percent} \enspace(\num{37}/\num{91}) & \SI{40.7}{\percent} \enspace(\num{37}/\num{91}) \\
\bottomrule
\end{tabular}
\end{table}

\begin{figure}[t]
\centering
\includegraphics[width=0.86\textwidth]{res_ablation_c.pdf}
\caption[Évaluation C : effet du correctif de frontière]{Couverture complète
inter-pages. Le correctif produit l'effet prédit par le diagnostic~: la
couverture du contrôleur passe de \num{27} à \num{36} questions sur \num{91}
à l'amorce par défaut, réduisant l'écart avec la référence de dix points à
un point.}
\label{fig-res_ablation_c}
\end{figure}

Après implémentation de l'expansion inter-documents
(\cref{sec-controleur_expansion}), la couverture à l'amorce par défaut passe
de \num{27} à \num{36} questions sur \num{91}, soit de \SI{29.7}{\percent} à
\SI{39.6}{\percent} --- exactement l'effet prédit par le diagnostic, mesuré
une seule fois après implémentation plutôt qu'ajusté empiriquement. Ce que le
correctif n'établit pas est une supériorité~: à l'amorce par défaut le
contrôleur reste derrière d'une question, et les intervalles de Wilson à
\SI{95}{\percent} des deux proportions se recouvrent presque intégralement
(\numrange{30.1}{49.8} pour le contrôleur, \numrange{31.1}{50.9} pour la
référence)~: les deux politiques sont statistiquement indiscernables sur cet
échantillon.

\subsection{Analyse}
\label{sec-res_c_analyse}

Le contrôleur devance la référence sur \num{7} questions et la suit sur
\num{8}~: un échange à peu près équilibré, non un déficit systématique. Il
ouvre en moyenne \num{4.86} nœuds contre \num{5.00} pour la référence, c'est-à-dire
qu'il élague parfois un nœud sous le seuil d'ouverture plutôt que de dépenser
son budget --- souhaitable en soi, mais sans gain de couverture. Douze
questions sur \num{91} obtiennent une couverture nulle~: le problème n'y est
pas la propagation de la frontière mais le classement lui-même, qui ne place
pas les bonnes tuiles en tête même sur la page de départ --- pointant une
seconde fois vers le scoreur, non vers la politique. Le correctif de
frontière traitait la question de \emph{où} la frontière peut aller, non
celle de \emph{à quel point elle juge bien la pertinence} une fois arrivée,
et le point d'extension nécessaire (\cref{sec-controleur_scoreur}) est déjà en
place.

% =====================================================================
\section{Évaluation D : le système assemblé}
\label{sec-res_ablation_d}
% =====================================================================

\subsection{Protocole et résultat}

Cette évaluation mesure le pipeline complet~: le contrôleur post-correctif
sélectionne les preuves (budget \num{5}, amorce \num{3}), le modèle
vision-langage lit chaque tuile ouverte, et un appel de synthèse final
produit la réponse. Elle porte sur un tirage aléatoire reproductible
(\code{random.Random(seed=0).sample}) de \num{20} des \num{91} questions
inter-pages de l'évaluation~C --- l'échantillon reste restreint en raison du
coût, environ deux minutes par tuile ouverte. Les réponses sont notées par
correspondance exacte et $F_1$ sur tokens après normalisation standard de
SQuAD~\citep{squad2016}.

\begin{table}[t]
\centering
\small
\caption[Évaluation D : précision bout-en-bout]{Précision bout-en-bout du
pipeline complet, \num{20} questions multi-sauts HotpotQA. La taille de
l'échantillon ne permet pas d'intervalle de confiance utile~; c'est le mode de défaillance du \cref{tab-res_erreurs} qui porte l'information.}
\label{tab-res_ablation_d}
\begin{tabular}{@{}l r@{}}
\toprule
Métrique & Valeur \\
\midrule
Correspondance exacte & \num{0.050} \enspace(\num{1}/\num{20}) \\
$F_1$ sur tokens & \num{0.057} \\
\bottomrule
\end{tabular}
\end{table}

Une réponse exactement correcte sur vingt, et un $F_1$ de \num{0.057}~: prise
seule, cette mesure indique que le système assemblé n'est pas prêt pour un
usage réel. Ce qui la rend informative est la \emph{nature} des erreurs.

\subsection{Taxonomie des erreurs et analyse}
\label{sec-res_d_erreurs}

\begin{table}[t]
\centering
\small
\caption[Taxonomie des erreurs bout-en-bout]{Répartition des \num{20}
prédictions du système assemblé, établie manuellement.}
\label{tab-res_erreurs}
\begin{tabularx}{\textwidth}{@{}l r X@{}}
\toprule
Catégorie & $n$ & Exemple de prédiction \\
\midrule
Abstention explicite & \num{14} & « Not specified in evidence », « Not found in evidence », « Not determinable from evidence » \\
Retournement oui/non erroné & \num{1} & « No » contre une référence « yes » \\
Partiellement correcte ($F_1 > 0$) & \num{1} & réponse hybride sur « ... Tenerife, Spain ... » contre « Canary Islands, Spain » ($F_1 = \num{0.133}$) \\
Entité erronée & \num{3} & « valuation » contre « business »~; « iMac G5 » contre « keyboard function keys » \\
Réponse exacte & \num{1} & « Blake Shelton » \\
\bottomrule
\end{tabularx}
\end{table}

\begin{figure}[t]
\centering
\includegraphics[width=0.88\textwidth]{res_ablation_d.pdf}
\caption[Mode de défaillance bout-en-bout]{Quatorze prédictions sur vingt
sont une abstention explicite~: le système reconnaît que les preuves
rassemblées ne suffisent pas, plutôt que de produire une réponse inventée.}
\label{fig-res_ablation_d}
\end{figure}

\subsection{Sous-rappel de preuves plutôt que fabulation}
\label{sec-res_d_analyse}

\textbf{Le mode dominant est l'abstention.} Quatorze prédictions sur vingt
déclarent explicitement que les preuves ne permettent pas de répondre~: le
lecteur ne réussit pas à interpréter les preuves qu'on lui donne, il constate
correctement qu'elles sont insuffisantes --- la défaillance est en amont,
dans la sélection. \textbf{Cette localisation est cohérente avec
l'évaluation~C}~: sa couverture complète de \SI{39.6}{\percent} signifie que
sur six questions sur dix, une question multi-sauts est structurellement
inrépondable et l'abstention y est la réponse correcte --- les deux
évaluations mesurent le même goulot par deux voies indépendantes.
\textbf{La fabulation est marginale}~: quatre prédictions sur vingt seulement
affirment quelque chose de faux, un comportement bien calibré face à des
preuves incomplètes, non garanti pour un système de RAG. \textbf{La synthèse
échoue aussi lorsque la couverture est complète}~: \SI{39.6}{\percent} fixe
un plancher optimiste que l'exactitude bout-en-bout (une réponse sur vingt)
n'atteint pas~; cet écart supplémentaire n'est pas encore diagnostiqué ---
preuve mal lue, invite sous-spécifiée, ou faiblesse de lecture du modèle de
base --- et reste à trier avant d'étendre l'échantillon
(\cref{sec-res_futur}). Le retournement oui/non mérite une mention à part~:
c'est la seule erreur susceptible d'induire un utilisateur en erreur sans
qu'il s'en aperçoive, un « non » assuré ne signalant rien de son incertitude.

% =====================================================================
\section{Coûts de calcul}
\label{sec-res_couts}
% =====================================================================

La génération du jeu de données coûte environ deux minutes par tuile en
exécution locale, soit de l'ordre de soixante heures de calcul séquentiel sur
\num{1784} tuiles~; l'émission concurrente des appels au sein d'une page
(\cref{sec-impl_concurrence}) et le cache sur disque (\cref{sec-impl_repro})
ramènent ce coût à un ordre praticable. L'entraînement contrastif de
l'évaluation~B a été mené sur une carte A100/H100 de \SI{80}{\giga\byte}
louée, lot de \num{4}, deux négatifs difficiles par question, trois époques
sur les \num{1339} exemples contrastifs~; le \anglais{gradient checkpointing}
(\cref{sec-contrastif_memoire}) coûte \SIrange{20}{30}{\percent} de temps
supplémentaire.

C'est le \textbf{coût du scoreur appris à l'inférence} qui explique la
principale expérience non menée~: substituer le lecteur affiné à TF-IDF dans
le contrôleur impose un passage avant du modèle par tuile candidate et par
question, là où un score lexical est essentiellement gratuit. Sur les
\num{91} questions de l'évaluation~C, chacune sur un graphe de plusieurs
dizaines de tuiles, cela représente plusieurs milliers de passages avant sur
GPU loué --- pour une seule configuration d'amorce. Ce coût, non une
difficulté d'implémentation, est ce qui a placé l'expérience hors de portée
du projet. À l'inverse, le rendu et le tuilage sont négligeables~: quelques
secondes par page, dominées par l'attente réseau, et la complexité
quadratique de l'élagage et du regroupement (\cref{sec-tuilage_regles}) ne se
manifeste jamais aux tailles de page réelles --- le tuilage guidé par le DOM
est en pratique gratuit, ce qui renforce le résultat de la \cref{sec-res_a_pixels}~: il ne s'agit pas d'échanger du temps de calcul
contre des tokens économisés, mais d'économiser les deux.

% =====================================================================
\section{Limites}
\label{sec-res_limites}
% =====================================================================

\paragraph{Les deux limites du contrôleur sont indépendantes.} Établies en
\cref{sec-res_synthese}. Le scoreur reste lexical alors que le lecteur affiné
de l'évaluation~B est une fonction de similarité bien meilleure, et la cause de
cette non-substitution est le coût de calcul (\cref{sec-res_couts}). La
politique, elle, n'emprunte que trois des sept relations instanciées
(\cref{tab-graphe_relations}), de sorte que la richesse du graphe reste
inexploitée quel que soit le scoreur employé. Les deux correctifs sont
distincts et peuvent être conduits séparément.

\paragraph{Confusion de famille de modèles.}
\label{sec-res_limites_famille}
Les questions synthétiques et le lecteur affiné proviennent de la même
famille Qwen --- une contrainte d'infrastructure (\cref{sec-contrastif_encodeur}),
Ollama n'exposant aucun accès aux poids pour l'affinage. Une partie du gain de
l'évaluation~B pourrait donc refléter les idiosyncrasies du générateur de
questions plutôt qu'une amélioration transférable à des questions humaines. Un
contrôle direct --- la même mesure de récupération sur les questions humaines
de HotpotQA --- est implémenté mais n'a pas été exécuté faute de budget de
calcul. Le seul indice disponible est la précision bout-en-bout de
\SI{5}{\percent} de l'évaluation~D, nettement inférieure aux chiffres de
l'évaluation~B et cohérente avec un déficit de généralisation, sans le
démontrer~: l'évaluation~D change aussi la tâche (pipeline complet, questions
inter-pages, mesure de réponse plutôt que de rang).

\paragraph{Le jeu contrastif ne couvre pas tout le corpus.}
\label{sec-res_limites_dataset}
Seize unités de page dont les paires question-réponse ont été générées ne
figurent pas dans le jeu contrastif, leur minage de négatifs n'ayant pas abouti
lors de la construction. Ces unités portent \num{378} des \num{1717} questions,
soit \SI{22}{\percent} du total pour \SI{8.6}{\percent} des unités. Elles sont
en moyenne deux fois et demie plus grandes que le reste du corpus, si bien que la
perte frappe préférentiellement les pages longues, celles-là mêmes que ce travail
vise. Le lecteur de l'évaluation~B est donc affiné sur un sous-ensemble du
corpus, et la proportion \num{80}/\num{20} annoncée en
\cref{sec-contrastif_split} porte sur ce sous-ensemble. Régénérer le minage sur
ces unités serait peu coûteux, le scoreur étant lexical, mais imposerait de
réaffiner l'adaptateur et de rejouer l'évaluation~B.

\paragraph{Portée limitée aux pages sémantiquement balisées.} Tous les
résultats quantitatifs portent sur Wikipédia, avec une liste de sélecteurs de
nettoyage propre à MediaWiki (\cref{sec-rendu})~; aucune vérification, même
qualitative, n'a été menée hors MediaWiki. La liste de sélecteurs est
spécifique~; la logique de tuilage --- boîtes englobantes, mesure serrée,
fusion, découpe, résolution des chevauchements --- ne l'est pas et ne connaît
rien de MediaWiki. L'approche n'est pas non plus revendiquée pour les
applications monopages hydratées par JavaScript ni les mises en page en
soupe de \code{div}, le régime que les agents web~\citep{webarena2023,
mind2web2023, seeact2024} étudient déjà côté action.

\paragraph{Le plafond de tuiles restreint la portée de l'évaluation~A.}
\label{sec-res_limites_plafond}
Le plafond de \num{25} tuiles (\cref{sec-res_a_protocole}) borne l'évaluation à
la portion supérieure des pages longues --- non biaisée vers du contenu
marginal, les infoboîtes et introductions y étant concentrées, mais partielle.
La valeur de \num{24.0} tuiles par page est un plafond atteint, non une mesure
de densité~: sans lui, le rapport d'environ cinq entre tuiles DOM et bandes de
grille serait plus élevé, dans une direction non mesurée pour le rapport de
pixels. Lever ce plafond exigerait de régénérer les paires question-réponse
sur l'ensemble du corpus, donc de répéter le coût de l'\cref{sec-res_couts}.

\paragraph{Tailles d'échantillon.} Dix-huit pages et \num{409} questions pour
l'évaluation~A. \num{330} questions de validation pour l'évaluation~B --- la
plus solide des quatre. \num{91} questions pour l'évaluation~C, dont les
intervalles de Wilson ne justifient qu'une indiscernabilité statistique.
\num{20} questions pour l'évaluation~D, qui ne supportent qu'une lecture
qualitative du mode de défaillance.

\paragraph{Autres limites.} La découpe des éléments démesurés
(\cref{sec-tuilage_regles}) reste indifférente au contenu à l'intérieur d'un
seul élément DOM --- une version étroite du défaut reproché à la grille. Le
parcours de corpus est peu profond (un saut, au plus trois tuiles porteuses de
liens, \cref{sec-graphe_interdoc})~; convient à HotpotQA, pas à une chaîne de
trois sauts ou davantage. Le minage de négatifs reste mono-page hors du chemin
HotpotQA. Tous les hyperparamètres de l'\cref{annexe-hyperparams} sont des
valeurs uniques justifiées mais non balayées, faute de budget de calcul. La
relation \code{semantic\_neighbor} et l'ablation de résolution
(\cref{sec-protocole_resolution}) sont implémentées mais non évaluées.

% =====================================================================
\section{Directions de travail}
\label{sec-res_futur}
% =====================================================================

Par ordre de bénéfice sur coût~: \textbf{(1)} brancher le lecteur affiné dans
le score de pertinence du contrôleur --- le point d'extension existe
(\cref{sec-controleur_scoreur}), seule l'exécution est coûteuse~; une
atténuation praticable serait un classement en deux étapes, TF-IDF pour
présélectionner, le lecteur affiné pour réordonner. \textbf{(2)} rejouer
l'évaluation~A avec ce même scoreur, pour trancher si l'avantage du texte
plat est un artefact du scoreur ou un résultat robuste. \textbf{(3)} exécuter
le contrôle de généralisation déjà écrit (\cref{sec-res_limites_famille}).
\textbf{(4)} faire exploiter au contrôleur \code{section\_hierarchy} et
\code{layout\_adjacency}, aujourd'hui jamais interrogées. \textbf{(5)} miner
les négatifs difficiles avec le lecteur lui-même une fois un premier tour
d'affinage effectué. \textbf{(6)} vérifier qualitativement le tuilage hors
MediaWiki, pour séparer ce qui relève du nettoyage de ce qui relève de la
logique de tuilage. \textbf{(7)} étendre l'évaluation~D aux
soixante-onze questions restantes de l'évaluation~C, une fois triée la cause
de l'écart entre couverture de preuves et exactitude bout-en-bout.

% =====================================================================
\chapter{Hyperparamètres}
\label{annexe-hyperparams}
% =====================================================================

Toutes les valeurs référencées dans le \cref{chap-methodologie} sont rassemblées
ici. Chacune est une valeur unique, documentée et centralisée dans un module de
configuration~; aucune n'a fait l'objet d'un balayage
(\cref{sec-res_limites}).

\begin{table}[H]
\centering
\small
\caption[Hyperparamètres : rendu, tuilage, graphe et scoreur lexical]{Hyperparamètres
de la partie hors ligne de la chaîne. Les valeurs en rouge restent à extraire du
module de configuration.}
\label{tab-annexe_hyperparams}
\begin{tabularx}{\textwidth}{@{}X l l@{}}
\toprule
\textbf{Paramètre} & \textbf{Valeur} & \textbf{Référence} \\
\midrule
\multicolumn{3}{@{}l}{\emph{Rendu et extraction}} \\
Largeur du \anglais{viewport} de rendu & \SI{2048}{\pixel} & \cref{sec-rendu} \\
Hauteur du \anglais{viewport} (sans effet, capture pleine page) & \SI{1000}{\pixel} & \cref{sec-rendu} \\
Facteur d'échelle de périphérique & \num{1} & \cref{sec-rendu} \\
Condition d'attente avant capture & \code{networkidle} & \cref{sec-rendu} \\
\addlinespace
\multicolumn{3}{@{}l}{\emph{Tuilage}} \\
Hauteur maximale de tuile (avant découpe) & \SI{2048}{\pixel} & \cref{sec-tuilage_regles} \\
Seuil de mise en tampon (hauteur d'élément) & \SI{80}{\pixel} & \cref{sec-tuilage_regles} \\
Tolérance d'imbrication (élagage) & \SI{1}{\pixel} & \cref{sec-tuilage_elagage} \\
Seuil de chevauchement $\rho_{\min}$ & \num{0.15} & \cref{sec-tuilage_chevauchement} \\
Hauteur de page avant découpage en sections & \SI{20480}{\pixel} \enspace($10 \times h_{\max}$) & \cref{sec-tuilage_sections} \\
Budget de hauteur d'une section trop haute & \aremplir[budget] & \cref{sec-tuilage_sections} \\
Tuiles retenues par unité de page (plafond) & \num{25} & \cref{sec-tuilage} \\
\addlinespace
\multicolumn{3}{@{}l}{\emph{Graphe}} \\
Tolérance de regroupement en rangées & \SI{20}{\pixel} & \cref{sec-graphe_ordre} \\
Tuiles porteuses de liens conservées ($k_{\text{liens}}$) & \num{3} & \cref{sec-graphe_interdoc} \\
Pages suivies depuis l'amorce ($m$) & \aremplir & \cref{sec-graphe_corpus} \\
Seuil de voisinage sémantique (cosinus) & \num{0.2} & \cref{sec-graphe_semantique} \\
Voisins sémantiques maximum par tuile & \num{3} & \cref{sec-graphe_semantique} \\
\addlinespace
\multicolumn{3}{@{}l}{\emph{Scoreur lexical TF-IDF}} \\
Plage de $n$-grammes & \aremplir & \cref{sec-controleur_scoreur} \\
Liste de mots vides & \aremplir & \cref{sec-controleur_scoreur} \\
Fréquences documentaires minimale et maximale & \aremplir & \cref{sec-controleur_scoreur} \\
Pondération sous-linéaire des fréquences & \aremplir[oui/non] & \cref{sec-controleur_scoreur} \\
Longueur de l'aperçu textuel par tuile & \aremplir[caractères] & \cref{sec-tuilage} \\
\bottomrule
\end{tabularx}
\end{table}

\begin{table}[H]
\centering
\small
\caption[Hyperparamètres : apprentissage, contrôleur, modèles et ablations]{Hyperparamètres
de l'entraînement, du contrôleur en ligne et des conditions de référence.}
\label{tab-annexe_hyperparams_b}
\begin{tabularx}{\textwidth}{@{}X l l@{}}
\toprule
\textbf{Paramètre} & \textbf{Valeur} & \textbf{Référence} \\
\midrule
\multicolumn{3}{@{}l}{\emph{Entraînement contrastif}} \\
Négatifs difficiles minés par question ($K$) & \num{2} & \cref{sec-contrastif_mining} \\
Rang LoRA ($r$) & \num{16} & \cref{sec-contrastif_lora} \\
Facteur d'échelle LoRA ($\alpha$) & \num{32} & \cref{sec-contrastif_lora} \\
\anglais{Dropout} LoRA & \num{0.05} & \cref{sec-contrastif_lora} \\
Paramètres entraînables & \aremplir & \cref{sec-contrastif_lora} \\
Optimiseur, taux d'apprentissage, échauffement & \aremplir & \cref{sec-contrastif_lora} \\
Température InfoNCE ($\tau$) & \num{0.05} & \cref{sec-contrastif_loss} \\
Taille de lot ($B$) & \num{4} & \cref{sec-contrastif_loss} \\
Époques & \num{3} & \cref{sec-res_couts} \\
Proportion du partitionnement (par article) & \num{80}/\num{20} & \cref{sec-contrastif_split} \\
\addlinespace
\multicolumn{3}{@{}l}{\emph{Contrôleur}} \\
Budget d'ouvertures ($B$) & \num{5} & \cref{sec-controleur_politique} \\
Amorce ($k_{\text{seed}}$) & \num{3} & \cref{sec-controleur_politique} \\
Seuil d'ouverture ($\theta_{\text{ouv}}$) & \num{0.05} & \cref{sec-controleur_politique} \\
Seuil d'activation ($\theta_{\text{act}}$) & \num{0.05} & \cref{sec-controleur_politique} \\
Appels au modèle concurrents & \num{4} & \cref{sec-impl_concurrence} \\
\addlinespace
\multicolumn{3}{@{}l}{\emph{Modèles et décodage}} \\
Modèle lecteur (inférence) & \code{qwen3-vl} (Ollama) & \cref{sec-implementation} \\
Modèle lecteur (affinage) & \code{Qwen2-VL-7B-Instruct} & \cref{sec-contrastif_encodeur} \\
Aire d'image minimale du processeur (\code{min\_pixels}) & \aremplir & \cref{sec-contrastif_memoire} \\
Aire d'image maximale du processeur (\code{max\_pixels}) & \aremplir & \cref{tab-cout_tokens} \\
Température de décodage & \num{0} \enspace(glouton) & \cref{sec-impl_repro} \\
Fenêtre de contexte & \num{8192} tokens & \cref{sec-impl_repro} \\
\addlinespace
\multicolumn{3}{@{}l}{\emph{Références et ablations}} \\
Seuil de découpage en sections, propre à l'évaluation A & \SI{51200}{\pixel} \enspace($25 \times h_{\max}$, contre $10\times$ pour le corpus principal) & \cref{sec-protocole_grille} \\
Hauteur de bande de la grille aveugle & \SI{1024}{\pixel} & \cref{sec-protocole_grille} \\
Taille de fenêtre du texte plat & \num{400} caractères & \cref{sec-protocole_grille} \\
Largeur de bande de la grille aveugle & \SI{2048}{\pixel} \enspace(largeur de rendu) & \cref{sec-protocole_grille} \\
Facteur d'échelle de compression (Lanczos) & \num{0.5} & \cref{sec-protocole_resolution} \\
Graine du tirage aléatoire, sous-échantillon évaluation D & \num{0} & \cref{sec-res_ablation_d} \\
\bottomrule
\end{tabularx}
\end{table}

% =====================================================================
\chapter{Invites utilisées}
\label{annexe-invites}
% =====================================================================

Les invites sont reproduites telles qu'elles figurent dans le module de
configuration. Celles adressées au modèle générateur de questions sont en
français, celles de la synthèse finale en anglais --- la langue des articles
Wikipédia évalués et des réponses de référence de HotpotQA.

\section{Génération de paires question-réponse}

Cette invite accompagne l'image d'une unique tuile
(\cref{sec-contrastif_qa}).

\begin{lstlisting}[language={}, frame=single, basicstyle=\ttfamily\scriptsize]
Tu vois un unique fragment recadré (une "tuile") d'une capture d'écran de
page web, pas la page complète.

Génère UNE question précise et autonome dont la réponse est visible dans
cette image, et UNIQUEMENT dans cette image (quelqu'un qui n'a pas vu cette
tuile ne pourrait pas deviner la réponse par hasard). Évite les questions
vagues ("de quoi parle cette image ?") : vise une question factuelle et
spécifique (un nombre, un nom, une date, une définition, une légende...).

Si l'image ne contient aucune information exploitable (décorative, vide,
illisible), réponds exactement: SKIP

Sinon, réponds STRICTEMENT en JSON, sans texte autour, au format :
{"question": "...", "answer": "..."}
\end{lstlisting}

Trois propriétés de cette invite portent le protocole. L'exigence
d'\emph{autonomie} de la question rend la paire utilisable comme exemple positif de
récupération~: une question dont la réponse serait devinable sans voir la tuile
n'apprendrait rien au modèle. La sortie \code{SKIP} explicite évite d'avoir à
interpréter un refus formulé librement. Et le format JSON strict, sans texte
d'encadrement, rend l'analyse syntaxique déterministe.

\section{Instruction d'encodage PromptEOL}

Cette instruction suit l'entrée --- question textuelle ou image de tuile,
indifféremment --- et l'état caché de la dernière couche à la position de son
dernier token est pris comme plongement (\cref{sec-contrastif_encodeur}).

\begin{lstlisting}[language={}, frame=single, basicstyle=\ttfamily\scriptsize]
\aremplir[instruction PromptEOL exacte, depuis le module de configuration]
\end{lstlisting}

Son rôle est de contraindre le modèle à comprimer le sens de l'entrée en une
position unique~; c'est son application \emph{à l'identique} aux deux modalités qui
place question et tuile dans le même espace, sans tête de projection additionnelle.

\section{Juge de faux négatifs}

Ce second filtre, optionnel et non activé pour les jeux de données évalués, est
adressé au modèle avec l'image d'une tuile candidate au statut de négatif
difficile (\cref{sec-contrastif_mining}).

\begin{lstlisting}[language={}, frame=single, basicstyle=\ttfamily\scriptsize]
Tu vois un fragment ("tuile") d'une capture d'écran de page web.

Question : {question}
Réponse attendue (trouvée dans une AUTRE tuile de la même page) : {answer}

Cette image permet-elle, à elle seule, de retrouver cette même information
-- même reformulée (synonyme), ou exprimée visuellement (couleur, logo,
forme) plutôt que textuellement ? Réponds STRICTEMENT par un seul mot :
OUI ou NON.
\end{lstlisting}

L'insistance sur les formulations visuelles est le point de cette invite~: c'est
exactement ce que le filtre lexical par sous-chaîne ne peut pas détecter.

\section{Lecture d'une tuile ouverte}

Adressée au modèle pour chaque tuile que le contrôleur a ouverte
(\cref{sec-controleur_lecture}).

\begin{lstlisting}[language={}, frame=single, basicstyle=\ttfamily\scriptsize]
Question : {query}
Extrait uniquement les informations de cette image utiles pour y répondre.
\end{lstlisting}

\section{Synthèse de la réponse finale}

Adressée au modèle en texte seul, une fois les preuves rassemblées.

\begin{lstlisting}[language={}, frame=single, basicstyle=\ttfamily\scriptsize]
Here are excerpts from a web page, selected as relevant evidence to answer
a question.

Evidence:
{evidence}

Question: {question}

Answer the question using ONLY this evidence, in English. If the evidence
does not allow you to answer, say so explicitly rather than making
something up.
\end{lstlisting}

La dernière phrase de cette invite est directement responsable du mode de
défaillance observé en \cref{sec-res_d_erreurs}~: elle instruit explicitement le
modèle de s'abstenir plutôt que d'inventer. Les quatorze abstentions sur vingt
prédictions sont donc un comportement conforme aux instructions face à des
preuves incomplètes, et non une défaillance du lecteur.

% =====================================================================
\chapter{Organisation du dépôt}
\label{annexe-depot}
% =====================================================================

Le code source est organisé avec un module par étage du pipeline. Cette
correspondance directe entre les sections du \cref{chap-methodologie} et les
modules est délibérée~: elle permet de localiser sans ambiguïté l'implémentation
de chaque décision de conception décrite.

\section{Modules du pipeline}

\begin{table}[H]
\centering
\small
\caption[Modules du pipeline]{Modules du pipeline et sections correspondantes.}
\label{tab-annexe_modules}
\begin{tabularx}{\textwidth}{@{}l X l@{}}
\toprule
\textbf{Module} & \textbf{Rôle} & \textbf{Section} \\
\midrule
\code{config.py} & constantes centralisées, invites & \cref{annexe-hyperparams} \\
\code{dom\_extraction.py} & rendu Playwright, extraction des boîtes englobantes, mesure serrée du texte coulant & \cref{sec-rendu} \\
\code{tiling.py} & élagage, regroupement des chevauchements, appariement, fusion, découpe & \cref{sec-tuilage} \\
\code{sectioning.py} & découpage des pages démesurées aux frontières \code{h1}/\code{h2} & \cref{sec-tuilage_sections} \\
\code{graph\_builder.py} & graphe à deux niveaux, sept relations, export XML & \cref{sec-graphe} \\
\code{corpus\_builder.py} & fusion multi-pages, préfixage, points d'entrée & \cref{sec-graphe_corpus} \\
\code{vlm\_client.py} & client vision-langage, cache sur disque, concurrence & \cref{sec-implementation} \\
\code{qa\_generation.py} & génération synthétique de paires question-réponse & \cref{sec-contrastif_qa} \\
\code{hard\_negative\_mining.py} & minage TF-IDF, double filtre anti-faux-négatif & \cref{sec-contrastif_mining} \\
\code{contrastive\_dataset.py} & jeu de données PyTorch, partitionnement par article & \cref{sec-contrastif_split} \\
\code{lora\_finetune.py} & tête PromptEOL, LoRA sur ViT et décodeur, perte InfoNCE & \cref{sec-contrastif} \\
\code{evidence\_controller.py} & machine à états, politique best-first, frontière inter-documents & \cref{sec-controleur} \\
\code{multihop\_qa.py} & construction des exemples inter-pages depuis HotpotQA & \cref{sec-protocole_multihop} \\
\code{image\_compression.py} & sous-échantillonnage Lanczos pour l'ablation de résolution & \cref{sec-protocole_resolution} \\
\code{visualization.py} & superpositions de tuiles, graphes interactifs & --- \\
\bottomrule
\end{tabularx}
\end{table}

\section{Scripts d'expérimentation}

\begin{table}[H]
\centering
\small
\caption[Scripts d'expérimentation]{Scripts d'expérimentation et évaluations
correspondantes.}
\label{tab-annexe_scripts}
\begin{tabularx}{\textwidth}{@{}l X l@{}}
\toprule
\textbf{Script} & \textbf{Rôle} & \textbf{Évaluation} \\
\midrule
\code{build\_dataset.py} & pilote le pipeline sur une liste de pages, sauvegarde par unité & \cref{sec-res_donnees} \\
\code{build\_multihop\_dataset.py} & construit les exemples HotpotQA, journalise les abandons & \cref{sec-res_donnees} \\
\code{grid\_baseline.py} & référence en grille aveugle, indépendante de \code{tiling.py} & A \\
\code{dom\_vs\_grid\_experiment.py} & construit les trois conditions et classe les questions & A \\
\code{eval\_retrieval.py} & récupération sur la partie de validation, avec et sans adaptateur & B \\
\code{eval\_retrieval\_tfidf.py} & référence lexicale sur le même partitionnement & B \\
\code{eval\_multihop\_controller.py} & couverture inter-pages, contrôleur contre top-$k$ & C \\
\code{eval\_multihop\_e2e.py} & pipeline complet, correspondance exacte et $F_1$ & D \\
\code{eval\_hotpotqa\_reader.py} & contrôle de généralisation \emph{(écrit, non exécuté)} & \cref{sec-res_limites_famille} \\
\code{compress\_dataset\_images.py} & produit le dossier miroir en résolution réduite & \cref{sec-protocole_resolution} \\
\code{generate\_paper\_figures.py} & superpositions de tuilage des \cref{fig-intro_teaser,fig-tuiles_reelles} & --- \\
\bottomrule
\end{tabularx}
\end{table}

\section{Tests}

{\sloppy
Les tests suivent la séparation par dépendances décrite en
\cref{sec-implementation}. Les tests de \code{tiling}, \code{graph\_builder},
\code{sectioning}, \code{hard\_negative\_mining}, \code{multihop\_qa},
\code{evidence\_controller}, \code{corpus\_builder} et
\code{image\_compression} s'exécutent sur des cas synthétiques, sans
réseau, sans modèle et sans GPU. Les tests de \code{dom\_extraction}
s'appuient sur Playwright et vérifient l'extraction de bout en bout contre une
page HTML fournie avec le dépôt, incluant un cas dédié à l'élément
flottant partageant sa rangée avec du texte
(\cref{sec-rendu_textflow}).
\par}

% =====================================================================
\chapter{Protocole de reproduction}
\label{annexe-repro}
% =====================================================================

\section{Installation}

\begin{lstlisting}[language=bash, frame=single, basicstyle=\ttfamily\scriptsize]
python3 -m venv .venv
source .venv/bin/activate
pip install -r requirements.txt
playwright install chromium

# Backend d'inférence (aucun GPU requis)
ollama pull qwen3-vl
\end{lstlisting}

\section{Vérification de l'installation}

Les tests sans dépendance lourde s'exécutent sur une page locale et ne
nécessitent ni réseau ni modèle.

\begin{lstlisting}[language=bash, frame=single, basicstyle=\ttfamily\scriptsize]
python3 -m pytest tests/ -x
python3 tests/test_pipeline_smoke.py
\end{lstlisting}

\section{Reproduction des évaluations}

\begin{lstlisting}[language=bash, frame=single, basicstyle=\ttfamily\scriptsize]
# Construction du jeu de données (long : ~2 min par tuile en local,
# mis en cache sur disque ; reprise automatique par unité de page)
python3 -m scripts.build_dataset
python3 -m scripts.build_multihop_dataset

# Évaluation A -- tuilage guidé, grille aveugle, texte plat
python3 -m scripts.dom_vs_grid_experiment

# Évaluation B -- lecteur affiné (nécessite un GPU pour l'entraînement)
python3 -m src.lora_finetune --epochs 3 --batch-size 4 \
        --gradient-checkpointing
python3 -m scripts.eval_retrieval
python3 -m scripts.eval_retrieval_tfidf

# Évaluation C -- contrôleur contre top-k statique
python3 -m scripts.eval_multihop_controller

# Évaluation D -- pipeline complet
python3 -m scripts.eval_multihop_e2e
\end{lstlisting}

\section{Application de démonstration}

\begin{lstlisting}[language=bash, frame=single, basicstyle=\ttfamily\scriptsize]
uvicorn webapp.backend.main:app --reload
# puis ouvrir http://127.0.0.1:8000/
\end{lstlisting}

Chaque requête recalcule l'intégralité du pipeline sur la page indiquée, sans
corpus préconstruit~: le temps de réponse va de quelques dizaines de secondes à
plusieurs minutes selon la page et le budget du contrôleur
(\cref{sec-implementation}).

\section{Notes de reproductibilité}

Trois points conditionnent la reproductibilité exacte des chiffres rapportés.
Le \textbf{décodage doit être glouton} (température nulle,
\cref{sec-impl_repro})~: une température positive rend les appels non
reproductibles, avec un symptôme trompeur --- la même question peut recevoir
une réponse correcte puis une abstention sans aucune différence dans les
tuiles ouvertes. La \textbf{fenêtre de contexte doit être portée à
\num{8192} tokens}~: à la valeur par défaut de \num{4096}, un raisonnement
long l'épuise avant la réponse, produisant un contenu vide silencieusement
interprété comme \code{SKIP}. Enfin, les résultats portent en grande partie
sur des \textbf{pages web vivantes}~; Wikipédia évolue, et les taux de
conservation des questions HotpotQA (\cref{sec-res_donnees}) dériveront avec
le temps. Pour l'évaluation~A, \code{scripts/pin\_wikipedia\_snapshot.py}
fige les révisions des dix-huit pages retenues, chargées automatiquement par
\code{dom\_vs\_grid\_experiment.py} avec repli sur le contenu en direct~; les
évaluations B, C et D n'en bénéficient pas encore.

\paragraph{Vérification de stabilité.} Une réexécution de l'évaluation~A sur
ce contenu figé (\cref{tab-annexe_repro_check}) reproduit le classement et
les conclusions qualitatives du \cref{tab-res_ablation_a} --- la grille garde
le meilleur R@1 mais le pire taux d'attachement, texte plat et tuilage DOM
restent proches et stables. Les écarts, de un à quatre points selon la
métrique, sont imputables à la dérive du contenu Wikipédia entre les deux
exécutions et ne changent aucune conclusion de la \cref{sec-res_ablation_a}.

\begin{table}[H]
\centering
\small
\caption[Vérification de stabilité de l'évaluation A]{Évaluation A, chiffres du
rapport (\cref{tab-res_ablation_a}, contenu Wikipédia en direct au moment du
\anglais{run}) contre une réexécution sur le contenu figé par
\code{pin\_wikipedia\_snapshot.py}. Mêmes dix-huit pages, mêmes
\num{409} questions~; le classement entre conditions et l'écart d'attachement
sont préservés.}
\label{tab-annexe_repro_check}
\begin{tabular}{@{}l r r r r@{}}
\toprule
& $n$ & R@1 & Attachement & \%rang \\
\midrule
\multicolumn{5}{@{}l}{\emph{Tuiles DOM}} \\
Rapport & \num{212} & \num{0.547} & \SI{51.8}{\percent} & \num{0.868} \\
Réexécution, contenu figé & \num{211} & \num{0.550} & \SI{51.6}{\percent} & \num{0.869} \\
\addlinespace
\multicolumn{5}{@{}l}{\emph{Grille aveugle}} \\
Rapport & \num{108} & \num{0.722} & \SI{26.4}{\percent} & \num{0.839} \\
Réexécution, contenu figé & \num{105} & \num{0.762} & \SI{25.7}{\percent} & \num{0.854} \\
\addlinespace
\multicolumn{5}{@{}l}{\emph{Texte plat}} \\
Rapport & \num{232} & \num{0.612} & \SI{56.7}{\percent} & \num{0.867} \\
Réexécution, contenu figé & \num{231} & \num{0.602} & \SI{56.5}{\percent} & \num{0.864} \\
\bottomrule
\end{tabular}
\end{table}

% =====================================================================
\chapter{Considérations éthiques}
\label{annexe-ethique}
% =====================================================================

\section{Accès aux données}

Le constructeur de corpus télécharge une page d'amorce et un nombre borné de
pages liées, avec temporisation et reprise plutôt que téléchargement agressif. Un
déploiement en production devrait néanmoins préférer l'API REST de Wikimedia ou un
\anglais{dump} statique au téléchargement en direct~; le prototype ne le fait pas
par défaut, et nous le signalons comme une limite de son mode de fonctionnement
plutôt que comme un choix défendable à grande échelle.

\section{Licences}

Le texte de Wikipédia est publié sous licence CC~BY-SA. Les images qui y sont
incorporées portent en revanche des licences hétérogènes, dont certaines non
libres. Les captures et les tuiles produites par le pipeline sont donc des œuvres
dérivées que nous ne redistribuons pas~; seul le code est publié.

\section{Statut des données synthétiques}

Les questions générées en \cref{sec-contrastif_qa} sont produites par un modèle et
n'ont pas été vérifiées par un humain. Elles doivent être étiquetées comme telles
en cas de diffusion. Cette réserve n'est pas formelle~: une question mal formée ou
dont la réponse est inexacte introduit du bruit dans l'entraînement contrastif, et
le taux de telles questions n'a pas été mesuré.

\section{Aptitude au déploiement}

Le mode de défaillance observé en \cref{sec-res_d_erreurs} est rapporté pour la
sécurité du déploiement autant que pour la complétude. L'erreur dominante est
l'abstention explicite, la plus sûre des deux pour un système de cette précision.
Mais une prédiction sur vingt était une réponse oui/non assurée et fausse
--- le mode de défaillance le plus susceptible d'induire en erreur un utilisateur
qui ne vérifierait pas une affirmation factuelle brève, puisqu'un « non » concis
ne signale rien de son incertitude. Un seul cas sur cet échantillon ne permet
pas d'estimer un taux, mais suffit à établir que le mode existe et n'est pas
détectable par l'utilisateur sans vérification externe.

Nous ne considérons donc pas le système apte au déploiement là où une réponse
oui/non erronée porterait des conséquences réelles, et nous rapportons ces
résultats précisément pour que cette décision ne se prenne pas sur la seule
motivation de la conception.
